\section{The mother problem and the first executed run}\label{sec:introduction}
A causal experiment supplies histories, not a manifold. Its geometry begins with the predictions of future experiments that can actually be performed. Even when finitely many such predictions form a sufficient state, a finite observer cannot retain an arbitrary point of that state space. Its actions determine the reports it will receive; its compression changes subsequent actions; and histories assigned the same label must be represented by one executable continuation. A risk bound must therefore concern the distribution acquired by that program, not only the set of formally reachable predictions.

A second distinction is equally important. An invariant calibrated observer is not automatically calibrated when started from the stated preparation. Converting an invariant assertion into a first-run assertion requires either an entrance estimate or a direct identification of the executed occupation law. The present article follows the second route. It proves a resource law for an independent geometric stopping experiment, whose mean number of acquisitions is finite. The stopping experiment is performed, rather than used as a formal abbreviation for an infinite-horizon cost.

The central construction has a dedicated initial label. The remaining labels are common cells of the predictive test image. Their conditional means are calibrated to the distribution at the actual stopping call. The calibration equation includes the preparation as an inhomogeneous term. Once the table has been fixed, its unique solution is the positive resolvent of the executed physical--observer instrument. Thus a fixed point in the construction is not an assertion that the machine starts in equilibrium. A discounted Bellman identity then cancels the continuation contribution and identifies the entire regret with one conditional Jensen gap under the actually acquired row mixture.

Under the experiment-level certificates in Definition~\ref{def:class}, Theorem~\ref{thm:main} gives
\[
 \calR_{\rm cp}(s,M;u_0)-V_s(u_0)
 \asymp \calR_{\rm hom}(s,M;u_0)-V_s(u_0)
 \asymp M^{-2/d},
\]
uniformly for $0<s\le1$ and $M\ge2$. Here $s$ is the prescribed independent stop probability, $1/s$ is the mean number of actual acquisition calls, and $V_s$ is the optimum over all full-history adaptive acquisition strategies in the same task. A completely autonomous implementation, including its halt status, uses at most twice the recurrent labels and satisfies the same order. The parameter $d$ is verified in the executable test image by actual acquired densities; it is not the dimension of an arbitrary parametrization of the experiment.

The exact calibration identity requires no feedback overlap. That certificate is used only to keep the Bellman oscillation, and hence the rate constant, uniform as the mean number of calls increases. No convergence or spectral gap of the finite observer is used. The table may depend on the preparation and on $s$: the assertion is uniform bounds, not one common table for all preparations and hazards. The finite-acquisition, uniform-audit and stationary results retained in the companion have different protocols and remain separate theorems \cite{GTFFinite}.

The raw realizations are controlled Gaussian observations and noncommuting quantum instruments. Both verify the same theorem before a code or budget is selected. A visible conditional expectation in the quantum instrument yields a changing acquired-rank extension. The lower retains the actual weight of every component. Geometric stopping is not arbitrary optimal stopping, and a finite mean budget is not a hard cap; Theorem~\ref{thm:cap} proves the capped comparison explicitly.

Conditional projection, discounted dynamic programming, positive resolvents and Kakutani's theorem are standard tools. The new assertion proved here is their preparation-injected causal calibration, the exact stopped regret identity, and its hard-state realization with matching acquired-mass lower bounds. Value-directed compression already relates belief approximation to decisions \cite{PB2002}; finite-memory control approximation and learning have a substantial literature \cite{KY2022,Kim2026,YB2008}. Section~\ref{sec:comparison} distinguishes their memory models and guarantees from this theorem. Neither the name of the program nor the preservation of earlier papers is a mathematical novelty claim.
